\documentclass[11pt,a4paper]{article}

\usepackage{amsmath,amssymb}
\usepackage{physics}
\usepackage{hyperref}
\usepackage[margin=2.5cm]{geometry}
\usepackage{booktabs}

\newcommand{\code}[1]{\texttt{#1}}
\newcommand{\Trsc}{\operatorname{Tr}_{sc}}

\title{\textbf{Disconnected qTMD/PDF One-Point Building Blocks}\\[4pt]
       \large Stochastic Quark Loops for Pion and Proton qTMD Analysis}
\author{PyQUDA Measurement Notes}
\date{\today}

\begin{document}
\maketitle
\tableofcontents

\section{Purpose}

This note documents the first PyQUDA implementation of disconnected qTMD/PDF
one-point quark loops.  The source file is
\code{pyquda\_measurement\_utils/Disconnected\_1pt\_qTMD\_vibe\_develop.py},
and the Perlmutter application is
\code{application/qTMD\_disconnected\_1pt/perlmutter}.

The loops are hadron independent.  Pion and proton information enters later,
when the loop is combined with the corresponding hadron two-point function in
ensemble analysis.

\section{Disconnected qTMD Loop}

For a stochastic source \(\xi_r\), define
\begin{equation}
  \eta_r = D^{-1}\xi_r,
\end{equation}
where \(D\) is the same clover Dirac operator used by the connected qTMD
measurements.  The implemented disconnected qTMD/PDF loop is
\begin{equation}
  L_{\Gamma,b}(q,\tau)
  =
  \frac{1}{N_r}
  \sum_{r=1}^{N_r}
  \sum_{\mathbf{x}}
  e^{i\mathbf{q}\cdot\mathbf{x}}\,
  \Trsc\!\left[
    \eta_r^\dagger(x)\,
    \Gamma\,
    O_b\,\xi_r(x)
  \right],
\end{equation}
where \(x=(\mathbf{x},\tau)\), \(\Gamma\) is one of the 16 bilinear gamma
matrices, and \(O_b\) is a nonlocal displacement or Wilson-line operator.

The current code convention is therefore
\begin{equation}
  \boxed{
  \mathrm{loop\ convention}
  =
  \eta^\dagger \Gamma O_b \xi
  }.
\end{equation}
The nonlocal operator \(O_b\) acts on the stochastic source side.  This choice
is intentionally minimal and makes the first implementation easy to compare
against connected PDF and local-current limits.

\section{Supported Operators}

The first implementation supports three operator kinds:
\begin{center}
\begin{tabular}{lll}
\toprule
\code{operator\_kind} & Operator & Status \\
\midrule
\code{GI\_PDF} & Straight-\(z\) gauge-invariant PDF Wilson line & Implemented \\
\code{CG\_PDF} & Straight-\(z\) coordinate-gauge displacement & Implemented \\
\code{CG\_qTMD} & qTMD-style coordinate-gauge displacement & Implemented \\
\bottomrule
\end{tabular}
\end{center}

\subsection{Straight-\(z\) PDF Operators}

For PDF operators,
\begin{equation}
  b = b_z \hat{z}.
\end{equation}
The coordinate-gauge version is a plain lattice shift:
\begin{equation}
  O_b^{\mathrm{CG\_PDF}}\xi(x) = \xi(x + b_z\hat{z}).
\end{equation}
The gauge-invariant version applies the straight Wilson line through repeated
covariant shifts:
\begin{equation}
  O_b^{\mathrm{GI\_PDF}}\xi(x)
  =
  U_z(x)U_z(x+\hat{z})\cdots
  \xi(x+b_z\hat{z}),
\end{equation}
with the corresponding backward-link product for negative \(b_z\).  In code
this is implemented by repeated calls to \code{gauge.pure\_gauge.covDev}.

\subsection{Coordinate-Gauge qTMD Operator}

The current qTMD operator follows the connected qTMD diagnostic convention:
\begin{equation}
  O_b^{\mathrm{CG\_qTMD}}\xi(x)
  =
  \xi(x + b_T \hat{e}_{\perp} + b_z\hat{z}),
\end{equation}
where \(\hat{e}_{\perp}\) is either \(\hat{x}\) or \(\hat{y}\).  The two
transverse directions are stored as \(\code{b\_X}\) and \(\code{b\_Y}\) in the
HDF5 output.

\paragraph{Current limitation.}
A fully gauge-invariant staple qTMD Wilson line is not implemented yet.  The
first production-quality disconnected qTMD extension should add the staple path
explicitly rather than relying on the coordinate-gauge displacement.

\section{Momentum Convention}

The loop momentum is the current momentum transfer,
\begin{equation}
  q = p_f - p_i.
\end{equation}
The one-point loop itself has no hadron source or sink momentum.  Those labels
enter when the loop is combined with a pion or proton two-point function.

\section{Combination with Hadron Two-Point Functions}

For a hadron two-point correlator \(C_2^H(p_f,t)\), the disconnected qTMD
three-point building block is formed as
\begin{equation}
  C^{H,\mathrm{disc}}_{3,\Gamma,b}
  (p_f,p_i;t,\tau)
  =
  \left\langle
    C_2^H(p_f,t)\,
    L_{\Gamma,b}(q,\tau)
  \right\rangle_{\mathrm{cfg}}
  -
  \left\langle C_2^H(p_f,t)\right\rangle_{\mathrm{cfg}}
  \left\langle L_{\Gamma,b}(q,\tau)\right\rangle_{\mathrm{cfg}},
\end{equation}
with
\begin{equation}
  q = p_f - p_i.
\end{equation}
The corresponding ratio is typically
\begin{equation}
  R^{H,\mathrm{disc}}_{\Gamma,b}(t,\tau)
  =
  \frac{
    C^{H,\mathrm{disc}}_{3,\Gamma,b}(p_f,p_i;t,\tau)
  }{
    C_2^H(p_f,t)
  },
\end{equation}
before applying the qTMD/PDF renormalization and soft-factor analysis.

The vacuum-subtraction term is essential because the one-point loop is hadron
independent.

\section{Stochastic Sources and Hierarchical Probing}

The source scheme is controlled by
\begin{equation}
  \code{noise\_scheme}
  \in
  \{\code{zn},\code{hierarchical\_probing}\}.
\end{equation}
For hierarchical probing, a base source \(\xi_b\) is multiplied by Rademacher
probing vectors \(h_k(x)=\pm 1\):
\begin{equation}
  \xi_{b,k}(x) = h_k(x)\xi_b(x).
\end{equation}
The number of effective inversions is
\begin{equation}
  N_{\mathrm{eff}}
  =
  N_{\mathrm{base}}\,
  N_{\mathrm{HP}}.
\end{equation}

The implemented HP ordering choices are
\begin{center}
\begin{tabular}{ll}
\toprule
Ordering & Meaning \\
\midrule
\code{global\_xyzt\_gray\_projected\_to\_evenodd}
& Gray code the full \(x,y,z,t\) site id. \\
\code{spatial\_xyz\_then\_t\_gray\_projected\_to\_evenodd}
& Gray code spatial sites first, then append time bits. \\
\bottomrule
\end{tabular}
\end{center}

Spin-color dilution and time dilution are not currently implemented.

\section{HDF5 Layout}

The output tag helper is
\code{get\_disconnected\_qTMD\_1pt\_file\_tag}.  The writer stores:
\begin{equation}
  \code{raw/loop\_pervec}
  \quad\mathrm{with\ shape}\quad
  [N_{\mathrm{eff}}, N_W, N_\Gamma, N_q, N_t].
\end{equation}
The raw bookkeeping datasets are
\begin{equation}
  \code{raw/source\_index},\quad
  \code{raw/base\_noise\_index},\quad
  \code{raw/hp\_index}.
\end{equation}

The averaged data are divided by the spatial volume and stored in a qTMD-like
layout:
\begin{equation}
  \code{avg/SS/<gamma>/PX...PY...PZ.../b\_X/eta0/bT.../bz...}.
\end{equation}
For \(y\)-direction transverse displacements the group is \(\code{b\_Y}\).

\section{Validation Status}

The first implementation has been checked with:
\begin{itemize}
  \item Python syntax compilation.
  \item A Perlmutter login-node smoke test for \(\code{GI\_PDF}\) at \(b_z=0\).
  \item A Perlmutter login-node smoke test for \(\code{GI\_PDF}\) at \(b_z=1\).
  \item A Perlmutter login-node smoke test for \(\code{CG\_qTMD}\) at
        \(b_T=1,b_z=1\).
\end{itemize}

\subsection{Local/PDF Limit Sanity Checks}

At zero displacement and zero momentum transfer, the different disconnected
operator implementations should collapse to the same local bilinear loop:
\begin{equation}
  O^{\mathrm{GI\_PDF}}_{b_z=0}
  =
  O^{\mathrm{CG\_PDF}}_{b_z=0}
  =
  O^{\mathrm{CG\_qTMD}}_{b_T=0,b_z=0}
  =
  1 .
\end{equation}
Therefore, for the same gauge field, stochastic source, gamma matrix, and
momentum \(q=0\), the expected equality is
\begin{equation}
  L^{\mathrm{GI\_PDF}}_{\Gamma,b_z=0}(0,\tau)
  =
  L^{\mathrm{CG\_PDF}}_{\Gamma,b_z=0}(0,\tau)
  =
  L^{\mathrm{CG\_qTMD}}_{\Gamma,b_T=0,b_z=0,\hat{x}}(0,\tau)
  =
  L^{\mathrm{CG\_qTMD}}_{\Gamma,b_T=0,b_z=0,\hat{y}}(0,\tau).
\end{equation}
This check is useful because it tests the HDF5 bookkeeping, Wilson-index
ordering, gamma loop ordering, and the coordinate-gauge/GI local limits without
requiring a hadron two-point function.

The current S8T32 Perlmutter login-node check used one stochastic vector with
matched random seed and found exact agreement:
\begin{equation}
  \max\left|
  L^{\mathrm{GI\_PDF}}
  -
  L^{\mathrm{CG\_PDF}}
  \right|
  =
  \max\left|
  L^{\mathrm{GI\_PDF}}
  -
  L^{\mathrm{CG\_qTMD}}_{\hat{x}}
  \right|
  =
  \max\left|
  L^{\mathrm{GI\_PDF}}
  -
  L^{\mathrm{CG\_qTMD}}_{\hat{y}}
  \right|
  =
  0 .
\end{equation}
The same equality was verified for both \(\code{raw/loop\_pervec}\) and the
averaged \(\Gamma=\gamma_5\) HDF5 path
\(\code{avg/SS/5/PX0PY0PZ0/b\_X/eta0/bT0/bz0}\), with the \(y\)-direction
branch stored under \(\code{b\_Y}\) for \(\code{CG\_qTMD}\).

The regression helper is
\code{tests/test\_qtmd\_disconnected\_local\_pdf\_limit.py}.  It skips
automatically if the local HDF5 smoke-test files are not present, and otherwise
requires exact equality for the local/PDF limit above.

\subsection{Nonzero-\(b_z\) Shift Sanity Check}

A second check probes the straight-\(z\) displacement convention at nonzero
separation.  For the coordinate-gauge operators, the PDF path at \(b_T=0\)
should be exactly the \(b_T=0\) limit of the qTMD coordinate shift:
\begin{equation}
  O^{\mathrm{CG\_PDF}}_{b_z=\pm 1}
  =
  O^{\mathrm{CG\_qTMD}}_{b_T=0,b_z=\pm 1,\hat{x}}
  =
  O^{\mathrm{CG\_qTMD}}_{b_T=0,b_z=\pm 1,\hat{y}} .
\end{equation}
This equality checks the sign of the \(z\)-shift, the Wilson-index ordering,
and the transverse-branch bookkeeping.  The corresponding S8T32 check found
exact agreement for both \(b_z=+1\) and \(b_z=-1\):
\begin{equation}
  \max\left|
  L^{\mathrm{CG\_PDF}}_{\Gamma,b_z=\pm1}
  -
  L^{\mathrm{CG\_qTMD}}_{\Gamma,b_T=0,b_z=\pm1,\hat{x}/\hat{y}}
  \right|
  =
  0 .
\end{equation}

The gauge-invariant PDF loop uses the same Wilson-index labels,
\([0,0,0,0]\), \([0,+1,0,0]\), and \([0,-1,0,0]\), but it is not expected to
equal the coordinate-gauge loop on a nontrivial gauge field because the
straight Wilson line is active:
\begin{equation}
  O^{\mathrm{GI\_PDF}}_{b_z=\pm1}
  =
  U_z \cdots \xi(x\pm\hat{z})
  \ne
  \xi(x\pm\hat{z})
  =
  O^{\mathrm{CG\_PDF}}_{b_z=\pm1}.
\end{equation}
The nonzero difference between \(\code{GI\_PDF}\) and \(\code{CG\_PDF}\)
therefore confirms that the GI path is not accidentally falling back to the
plain coordinate shift.

The regression helper is
\code{tests/test\_qtmd\_disconnected\_nonzero\_bz.py}.  During this check it
also exposed and fixed a transverse-branch bug in the \(\code{CG\_qTMD}\)
operator: when switching from the \(\hat{x}\) branch to the \(\hat{y}\) branch,
the shifted stochastic source must be reset to the local source before applying
the new branch displacement.

The next physics validation should compare the local/PDF limits against the
connected pion and proton qTMD workflows after constructing the disconnected
ratio with matched two-point functions.

\end{document}
